% !TEX root = ../main.tex
\chapter{変更を支える設計}
\label{chap:software-engineering}

\begin{chapterabstract}
実行中の一要求だけを見れば、プログラムは命令とデータの集まりです。しかし、そのAIアプリは翌日にも変更されます。配送会社が増え、注文データの置き場所が変わり、モデルの呼び出し方が更新されます。多くの開発者が触れるコードでは、動くことに加えて、変え続けられることが必要です。

ここでは時間を少し巻き戻し、この要求を処理するアプリがどう設計されたかを見ます。注文の事実、HTTPの形式、データベース、AIによる文章化を同じ場所へ混ぜず、変化の理由ごとに境界を引きます。抽象化、情報隠蔽、オブジェクト指向、関数型、テストは別々の流派ではなく、変更の波を小さくするための道具として登場します。
\end{chapterabstract}


\section{ソフトウェア危機}
\label{sec:04-software-engineering-01}

一件の要求を動かすだけなら、必要な処理を一つの場所へ書くこともできます。難しさが表面化するのは、利用者、機能、開発者、運用年数が増えたときです。まず、規模の拡大がソフトウェアへ何を起こしてきたかを歴史から見ます。

\begingroup
\begingroup
\paragraph{1968年NATO会議}
        大規模ソフトウェアの開発、配布、保守に共通する問題を議論するため、NATO科学委員会が会議を開催しました。
        翌年の報告書は、議論を「Software Engineering」の題で記録しています。

\paragraph{問題の構図}
        要求や環境が変わる一方、規模が増すほど全体の理解、見積もり、統合、検証、保守が難しくなります。

\par
\endgroup
\medskip
\begingroup
\paragraph{工学的な開発}
        言語やアルゴリズムだけでなく、設計、工程、検証、協働、運用を含む仕組みが必要です。

\paragraph{現在への読み替え}
        危機は一度で解消した事件ではありません。依存関係と変更速度を管理する課題として、現代の開発にも残ります。

\par
\endgroup
  \source{NATO Science Committee, \textit{Software Engineering: Report on a Conference}, 1969}
\par
\endgroup


\section{規模と相互作用}
\label{sec:04-software-engineering-02}

コード量が二倍になれば難しさも単純に二倍になるわけではありません。部分同士の関係が増え、ある変更が思いがけない場所へ伝わるからです。注文データ、配送規則、HTTP、LLMが互いを直接知るほど、送信処理の全体像はつかみにくくなります。

\begingroup
\begingroup
      \centering
      \begin{tikzpicture}[every node/.append style={font=\scriptsize}]
        \foreach \a in {90,162,234,306,378}{
          \node[node] (n\a) at (\a:15mm) {};
        }
        \foreach \a/\b in {90/162,90/234,90/306,90/378,162/234,162/306,162/378,234/306,234/378,306/378}{
          \draw[draw=RuleGray] (n\a)--(n\b);
        }
        \node[font=\small,align=center] at (0,-23mm) {5要素を全結合すると\\関係は10本};
      \end{tikzpicture}
\par
\endgroup
\medskip
\begingroup
\paragraph{本質的な複雑さ}
        問題領域に元からある規則、例外、状態、同時実行、外部連携です。単純化しても消せない部分があります。

\paragraph{偶発的な複雑さ}
        重複、曖昧な責務、循環依存、過剰な抽象化など、実装や開発方法が持ち込む難しさです。

\par
\endgroup
  \vspace{3mm}
  \takeaway{変更コストは行数だけでなく、理解して確認すべき依存関係の広さで増えます。}
\par
\endgroup


\section{分割統治}
\label{sec:04-software-engineering-03}

全体を一度に理解できないなら、責任を小さな部分へ分けます。ただし、切り分ければ終わりではありません。各部分が解く問いを決め、必要な情報だけを境界越しに渡し、最後に一つの応答へ組み戻せる分け方が必要です。

\begingroup
\begingroup
\paragraph{分解}
        入力、計算、保存、表示などの責務を見分け、名前を付けます。単にファイル数を増やすことではありません。

\par
\endgroup
\medskip
\begingroup
\paragraph{関心の分離}
        ある判断をするとき、無関係な詳細を同時に考えずに済む境界を作ります。

\par
\endgroup
\medskip
\begingroup
\paragraph{再統合}
        各部分の契約を定め、組み合わせたときのデータと制御の流れを検証します。

\par
\endgroup
  \vspace{5mm}
\paragraph{良い分割を評価する問い}
    一つの要求変更はどの部分へ届くか、部分を単独で理解・検証できるか、境界を越える情報は必要最小限かを確認します。

\par
\endgroup


\section{凝集度と結合度}
\label{sec:04-software-engineering-04}

良い分割かどうかは、ファイル数やクラス数では決まりません。同じ理由で変わる処理が一か所へ集まり、異なる理由で変わる部分が必要以上に結び付いていないかを見ます。その二つの向きを、凝集度と結合度で捉えます。

\begingroup
\begingroup
\paragraph{凝集度}
        一つのモジュール内の要素が、まとまった責務へどれだけ集中しているかを表す観点です。
        「なぜ一緒にあるか」を一文で説明できる状態を目指します。

\paragraph{低凝集の兆候}
        「共通」「便利」という理由だけで無関係な処理を集めると、利用者も変更理由も増えます。

\par
\endgroup
\medskip
\begingroup
\paragraph{結合度}
        モジュール間の依存の強さを表す観点です。共有する可変状態、内部構造への依存、呼び出し順の暗黙条件は結合を強めます。

\paragraph{低結合の手段}
        小さなインタフェース、不変なデータ、明示した契約、依存方向の統一で、境界越しの知識を減らします。

\par
\endgroup
  \takeaway{凝集度と結合度は数値だけで決めず、想定する変更に対して評価します。}
\par
\endgroup


\section{情報隠蔽}
\label{sec:04-software-engineering-05}

境界を引いても、内部のデータ形式や外部APIの癖がそのまま漏れれば、変更は境界を越えて広がります。変わりやすい判断を内側へ隠し、外からは安定した約束だけを見せることで、波及する範囲を抑えます。

\begingroup
  \centering
  \begin{tikzpicture}[node distance=4mm,every node/.append style={font=\scriptsize}]
    \node[box] (client1) {利用者A};
    \node[box,below=of client1] (client2) {利用者B};
    \node[accentbox,right=15mm of $(client1)!0.5!(client2)$] (api) {安定した\\インタフェース};
    \node[box,right=12mm of api] (detail) {変更し得る詳細\\保存形式・アルゴリズム};
    \draw[flow] (client1)--(api);
    \draw[flow] (client2)--(api);
    \draw[flow] (api)--(detail);
  \end{tikzpicture}
  \vspace{3mm}
\begingroup
\paragraph{隠す対象}
        データ表現、選択したアルゴリズム、外部サービス、キャッシュ方針など、将来変わり得る判断です。

\par
\endgroup
\medskip
\begingroup
\paragraph{得られる効果}
        判断を変更しても、公開契約を守る限り利用側の修正を減らせます。

\par
\endgroup
\medskip
\begingroup
\paragraph{分解基準}
        処理段階だけでなく、同時に変わる設計判断のまとまりでも境界を検討します。

\par
\endgroup
  \source{D. L. Parnas, “On the Criteria To Be Used in Decomposing Systems into Modules”, 1972}
\par
\endgroup


\section{構造化}
\label{sec:04-software-engineering-06}

部品同士の境界だけでなく、一つの部品の中にも読みやすい流れが必要です。制御が任意の場所へ飛ぶと、入力がどの経路を通って応答になるかを追えません。処理の進み方を限られた形へ整え、局所的に理由を説明できるようにします。

\begingroup
\begingroup
\paragraph{goto依存の問題}
        任意の位置への分岐が増えると、現在地点へ至る経路と、成立している状態を局所的に説明しにくくなります。

\paragraph{構造化の目的}
        問題は命令名ではなく、制御フローと状態変化を人が対応づけられるかです。例外処理や早期returnも同じ観点で読みます。

\par
\endgroup
\medskip
\begingroup
      \centering
      \begin{tikzpicture}[node distance=4mm,every node/.append style={font=\scriptsize}]
        \node[accentbox] (seq) {逐次\\Aの後にB};
        \node[accentbox,right=of seq] (sel) {選択\\条件で分岐};
        \node[accentbox,right=of sel] (loop) {反復\\条件まで繰返し};
        \draw[flow] (seq)--(sel);
        \draw[flow] (sel)--(loop);
      \end{tikzpicture}
      \vspace{5mm}
\paragraph{基本となる制御構造}
        逐次、選択、反復を明示的に組み合わせ、まとまりごとに事前条件・事後条件を考えます。

\par
\endgroup
  \source{E. W. Dijkstra, “Go To Statement Considered Harmful”, 1968}
\par
\endgroup


\section{抽象化}
\label{sec:04-software-engineering-07}

構造を整えても、毎回すべての細部を読んでいては大きなシステムを扱えません。注文を「行の集合」ではなく業務上の注文として、モデル呼び出しを「HTTPの手順」ではなく応答生成として扱うため、必要な性質だけを表面へ出します。

\begingroup
\begingroup
\paragraph{データ抽象}
        「残高」を整数の置き場所ではなく、入金・出金・参照という許された操作で捉えます。

\par
\endgroup
\medskip
\begingroup
\paragraph{手続き抽象}
        処理の実装手順を関数名と入出力の背後へ隠し、呼び出し側は目的を記述します。

\par
\endgroup
\medskip
\begingroup
\paragraph{制御抽象}
        イテレータや高階関数で、走査順や反復の詳細を再利用可能な形へまとめます。

\par
\endgroup
  \vspace{5mm}
\paragraph{抽象化の粒度}
    詳細を隠し過ぎると性能や失敗条件が見えず、隠さな過ぎると利用側へ変更理由が漏れます。利用者が正しく判断するための情報は契約に残します。

\par
\endgroup


\section{インタフェース}
\label{sec:04-software-engineering-08}

抽象化した部品を協調させるには、呼び出せる操作だけでなく、受け入れる条件、成功後の状態、失敗の表し方を決める必要があります。その約束が曖昧なら、要求が境界を越えるたびに意味がずれます。

\begingroup
\begingroup
\paragraph{事前条件}
        呼び出し側が満たす入力範囲、状態、権限などです。例：金額は0より大きい。

\paragraph{事後条件}
        正常終了時に提供側が保証する結果と状態変化です。例：保存成功なら注文IDを返す。

\par
\endgroup
\medskip
\begingroup
\paragraph{不変条件}
        操作の前後で保つべき性質です。例：在庫数は負にならない。

\paragraph{失敗の契約}
        タイムアウト、重複、入力不正、部分失敗を、戻り値・例外・再試行可能性として区別します。

\par
\endgroup
  \vspace{4mm}
  \takeaway{型シグネチャだけでは、順序、単位、境界値、副作用、失敗時の状態まで表せない場合があります。}
\par
\endgroup


\section{モデル}
\label{sec:04-software-engineering-09}

コード全体をそのまま見せても、構造を知りたい人と時間順を知りたい人では必要な情報が違います。ある問いに答える要素だけを選び、ほかを省いた表現がモデルです。省略するからこそ、重要な関係が見えるようになります。

\begingroup
\begingroup
\paragraph{静的構造}
        構成要素、属性、関連、依存など、時間を止めたときの関係を表します。クラス図やコンポーネント図が使えます。

\paragraph{答えたい問い}
        誰がどの情報を所有するか、どの方向へ依存するか、交換可能な境界はどこかを確認します。

\par
\endgroup
\medskip
\begingroup
\paragraph{動的な振る舞い}
        時間に沿うメッセージ、状態遷移、分岐、並行処理を表します。シーケンス図や状態機械図が使えます。

\paragraph{モデルの限界}
        図は実装そのものではありません。図の目的、対象範囲、抽象度、更新責任を明記し、コードやテストと矛盾させません。

\par
\endgroup
  \source{Object Management Group, \textit{UML 2.5.1 Specification}}
\par
\endgroup


\section{クラス図}
\label{sec:04-software-engineering-10}

まず、実行を止めた状態でアプリの登場人物を並べます。注文、明細、配送予定といった型がどんな情報を持ち、互いにどう結び付くかを示すと、データと責務の静的な地図ができます。

\begingroup
\begingroup
      \centering
      \begin{tikzpicture}[every node/.append style={font=\scriptsize},node distance=13mm]
        \node[box,align=left] (order) {\textbf{Order}\\\rule{24mm}{.3pt}\\id: OrderID\\items: []Item\\\rule{24mm}{.3pt}\\Total(): Money};
        \node[box,align=left,right=of order] (item) {\textbf{Item}\\\rule{23mm}{.3pt}\\qty: int\\price: Money\\\rule{23mm}{.3pt}\\Subtotal(): Money};
        \node[box,align=left,below=10mm of order] (repo) {\textbf{\textless\textless interface\textgreater\textgreater}\\OrderRepository\\\rule{29mm}{.3pt}\\Save(Order): error};
        \draw[draw=MutedBlue] (order)--node[above,font=\tiny]{1\quad contains\quad 1..*}(item);
        \draw[flow] (order)--(repo);
      \end{tikzpicture}
\par
\endgroup
\medskip
\begingroup
\paragraph{読み方}
        箱は型、区画は属性と操作、線は関連や依存を表します。多重度は関係できる個数です。

\paragraph{載せる情報}
        設計上の問いに必要な要素だけを載せます。全フィールドと全メソッドの転記を目的にしません。

\par
\endgroup
\par
\endgroup


\section{シーケンス図}
\label{sec:04-software-engineering-11}

静的な地図だけでは、利用者の質問が誰から誰へ渡るかは分かりません。時間を上から下へ流し、ブラウザ、アプリ、注文リポジトリ、配送サービス、モデルのやり取りを並べると、一要求の物語が設計図になります。

\begingroup
  \centering
  \begin{tikzpicture}[x=0.9mm,y=0.52mm,every node/.append style={font=\scriptsize}]
    \foreach \x/\name in {10/利用者,42/API,76/注文サービス,112/DB}{
      \node[box,minimum width=22mm] (h\x) at (\x,0) {\name};
      \draw[dashed,draw=PaleGray] (\x,-5)--(\x,-53);
    }
    \draw[flow] (10,-10)--node[above]{POST /orders}(42,-10);
    \draw[flow] (42,-20)--node[above]{Create(cmd)}(76,-20);
    \draw[flow] (76,-30)--node[above]{Save(order)}(112,-30);
    \draw[softflow] (112,-39)--node[above]{id}(76,-39);
    \draw[softflow] (76,-47)--node[above]{結果}(42,-47);
    \node[font=\tiny,align=left] at (112,-58) {上から下へ時間が進みます};
  \end{tikzpicture}
  \vspace{1mm}
\begingroup
\paragraph{見えるもの}
        呼び出し順、同期・非同期、戻り値、代替経路を表せます。

\par
\endgroup
\medskip
\begingroup
\paragraph{設計での用途}
        責務の偏り、往復回数、失敗時の戻り先、処理境界を検討します。

\par
\endgroup
\par
\endgroup


\section{オブジェクト指向}
\label{sec:04-software-engineering-12}

地図に現れた責務をコードへ移す一つの方法が、状態と操作をオブジェクトとしてまとめる考え方です。注文自身が不変条件を守れば、どの呼び出し元も同じ業務規則を使えます。重要なのは名詞をクラスへ変えることではなく、判断の持ち主を定めることです。

\begingroup
\begingroup
\paragraph{オブジェクト}
        同一性、状態、操作を持つ概念として扱います。設計では「何を知り、何を行う責任があるか」を考えます。

\paragraph{責務の配置}
        値を所有する対象の近くへ、その値の不変条件を守る操作を置くと、規則の重複を減らせます。

\par
\endgroup
\medskip
\begingroup
\paragraph{カプセル化}
        内部状態への任意の変更を許さず、公開操作を通じて不変条件を守ります。フィールドを非公開にするだけでは完了しません。

\paragraph{注意点}
        何でも一つの「Manager」へ集めると責務が肥大化します。データと規則が変わる理由を手掛かりに分けます。

\par
\endgroup
\par
\endgroup


\section{継承とコンポジション}
\label{sec:04-software-engineering-13}

既存の振る舞いを再利用するとき、型の上下関係として受け継ぐ方法と、必要な部品を組み合わせる方法があります。配送会社やモデル提供者のように差し替わる相手では、どちらを選ぶかが将来の変更範囲を大きく左右します。

\begingroup
\begingroup
\paragraph{継承}
        派生型を基底型として置換できる関係を表します。実装再利用だけを目的にすると、基底型の変更や前提へ強く結合します。

\paragraph{使う条件}
        公開契約と不変条件を派生型でも守り、利用者が型を置き換えても期待が崩れないかを確認します。

\par
\endgroup
\medskip
\begingroup
\paragraph{合成}
        対象が別の部品を持ち、委譲して機能を組み立てます。部品をインタフェース越しに交換しやすくできます。

\paragraph{判断例}
        「通知はメールである」ではなく「通知サービスが送信手段を持つ」と捉え、メールとSMSを交換可能な部品にします。

\par
\endgroup
  \takeaway{継承と合成は排他的な規則ではありません。置換関係と変更理由に合う方を選びます。}
\par
\endgroup


\section{ポリモーフィズム}
\label{sec:04-software-engineering-14}

異なる配送サービスが同じ「到着予定を調べる」という約束を満たせば、利用側は個々のAPI形式を知る必要がありません。同じ呼び方の向こうで実装を差し替えられる性質が、外部の変化を境界の外へ押し戻します。

\begingroup
\begingroup
      \begin{lstlisting}[language=Go,basicstyle=\ttfamily\fontsize{7}{8.3}\selectfont]
type Notifier interface {
    Send(ctx context.Context, m Message) error
}

type Checkout struct {
    notifier Notifier
}

func (c Checkout) Complete(ctx context.Context) error {
    // 注文処理
    return c.notifier.Send(ctx, Message{/* ... */})
}
      \end{lstlisting}
\par
\endgroup
\medskip
\begingroup
\paragraph{同じ呼び方、異なる実装}
        \code{Notifier}の契約を満たすメール、SMS、テスト用偽物を、呼び出し側を変えずに差し替えられます。

\paragraph{依存方向と粒度}
        上位の業務処理が必要な抽象を定義し、外部連携がそれを実装すると、業務処理から技術詳細への直接依存を減らせます。
        利用側が必要とする操作だけをインタフェースへ含めます。

\par
\endgroup
\par
\endgroup


\section{関数型の考え方}
\label{sec:04-software-engineering-15}

オブジェクトへ責務を置いても、状態の変化が多すぎれば一要求の結果を再現しにくくなります。入力から出力を求める純粋な計算と、DBやネットワークへ触れる副作用を分けると、注文の事実を応答材料へ変える部分を安定して検証できます。

\begingroup
\begingroup
\paragraph{純粋関数}
        同じ引数へ同じ結果を返し、外部から観測できる状態を変更しません。局所的な推論とテストをしやすくします。

\paragraph{副作用}
        入出力、時刻取得、乱数、共有状態の更新などです。必要な処理ですが、境界へ集めると計算部分を分離できます。

\par
\endgroup
\medskip
\begingroup
\paragraph{不変性}
        既存値を変更せず、新しい値を作ります。ある時点の値が後から変わらないため、共有や並行処理の推論を単純にできます。

\paragraph{対価}
        新しい値の確保やコピーが増える場合があります。永続データ構造、最適化、局所的な可変性を含めて測定します。

\par
\endgroup
\par
\endgroup


\section{高階関数}
\label{sec:04-software-engineering-16}

計算そのものを値として渡せると、繰り返し、絞り込み、変換といった流れを共通化できます。会話履歴から必要な発言を選び、モデルへ渡す形へ変える処理も、制御の骨格と個々の判断を分けて表せます。

\begingroup
\begingroup
      \begin{lstlisting}[language=Go]
prices := []int{1200, 800, 2500, 300}

// filter: 条件を満たす値だけ残す
eligible := Filter(prices, func(x int) bool {
    return x >= 1000
})

// map: 各値を別の値へ写す
taxed := Map(eligible, func(x int) int {
    return x * 110 / 100
})

// reduce: 列を一つの値へ畳み込む
total := Reduce(taxed, 0, func(a, x int) int {
    return a + x
})
      \end{lstlisting}
\par
\endgroup
\medskip
\begingroup
\paragraph{高階関数}
        関数を引数に受け取る、または関数を返す関数です。走査、選別、変換、集約の共通形を再利用します。

\paragraph{順序・副作用・Go}
        評価順序に依存する副作用をコールバックへ入れると、置換や並列化が難しくなります。
        Goは第一級関数を持ちます。上の\code{Filter}などは説明用の自作関数です。

\par
\endgroup
\par
\endgroup


\section{パラダイムの選択}
\label{sec:04-software-engineering-17}

ここまでの道具は、互いに排他的な陣営ではありません。寿命を持つ注文にはオブジェクト、値の変換には純粋関数、全体の手順には構造化された制御というように、扱う問題へ合わせて組み合わせます。

\begingroup
  \noindent
  \begin{tabularx}{\textwidth}{@{}lYYY@{}}
    \toprule
    観点 & 手続き型 & オブジェクト指向 & 関数型の考え方\\
    \midrule
    中心 & 操作手順と状態変化 & 責務を持つ対象の協調 & 値の変換と関数合成\\
    適する例 & 順序が本質の処理、装置制御 & 状態と規則を持つ業務概念 & 集計、変換、並列化しやすい計算\\
    注意 & 共有状態が広がる & 過剰な階層、肥大した対象 & 作用の扱い、割り当てコスト\\
    \bottomrule
  \end{tabularx}
  \vspace{4mm}
\paragraph{パラダイムの併用}
    Goでも、手続き的な制御、構造体とインタフェース、関数値と不変な値を同じプログラムで使えます。
    ラベルで統一するより、局所的な問題へ合う制約を選びます。

\par
\endgroup


\section{テスト}
\label{sec:04-software-engineering-18}

境界を設計しただけでは、その約束が守られる保証はありません。注文番号が存在しない場合、配送サービスが遅い場合、モデルが失敗した場合を、部品単独と接続後の両方で確かめます。テストは実装の検査であると同時に、境界の具体例です。

\begingroup
\begingroup
\paragraph{単体テスト}
        関数や小さなモジュールを隔離し、入力と結果、不変条件、境界値、失敗を高速に確認します。

\paragraph{設計へのフィードバック}
        外部状態なしに試せない、準備が大きい、結果を観測できない場合は、責務や依存境界を見直す手掛かりになります。

\par
\endgroup
\medskip
\begingroup
\paragraph{結合テスト}
        DB、ネットワーク、ファイル、複数モジュールの契約が、実際の組み合わせで成立するかを確認します。

\paragraph{二つのテストの役割}
        単体テストだけでは設定や通信形式の誤りを見逃し、結合テストだけでは原因特定と網羅が難しくなります。

\par
\endgroup
  \takeaway{テストの数ではなく、どのリスクをどの層で検出するかを設計します。}
\par
\endgroup


\section{品質の維持}
\label{sec:04-software-engineering-19}

テストが通った時点で設計が完成するわけでもありません。要求が加わるたびに構造は少しずつ崩れ得ます。小さく整え、他者の視点で読み、先送りした判断を見える形へ残すことで、次の変更が安全に通れる道を保ちます。

\begingroup
\begingroup
\paragraph{リファクタリング}
        外部から観測できる振る舞いを保ちながら、内部構造を改善します。テストで振る舞いの保持を確認します。

\par
\endgroup
\medskip
\begingroup
\paragraph{コードレビュー}
        正しさ、読みやすさ、保守性、セキュリティ、テストを別の視点で確認し、知識を共有します。

\par
\endgroup
\medskip
\begingroup
\paragraph{技術的負債}
        短期の便益のため将来の変更コストを増やす判断です。意図、利息、返済条件を可視化します。

\par
\endgroup
  \vspace{5mm}
\paragraph{変更の分離}
    振る舞いを変える修正と構造だけを変える修正を小さく分けると、レビューで確認すべき差と障害時の切り分けが明確になります。

  \source{Martin Fowler, \textit{Refactoring}; Ward Cunningham, “The WyCash Portfolio Management System”}
\par
\endgroup


\section{ケース：注文確認AI}
\label{sec:04-software-engineering-20}

ここで冒頭の要求を設計図へ戻します。AIが配送予定を勝手に推測するのではなく、注文と配送の仕組みが事実を返し、モデルはその事実を利用者向けの文へ整えます。この責任分担を、内側の業務規則と外側の技術へ分けます。

\begingroup
  \centering
  \begin{tikzpicture}[node distance=2mm,every node/.append style={font=\scriptsize}]
    \node[box,minimum width=92mm,minimum height=6mm,inner ysep=.5mm] (adapter1) {入力アダプタ：HTTP API／チャット画面};
    \node[accentbox,minimum width=92mm,minimum height=6mm,inner ysep=.5mm,below=of adapter1] (usecase) {ユースケース：質問受付・事実取得・応答生成};
    \node[box,minimum width=92mm,minimum height=6mm,inner ysep=.5mm,below=of usecase] (domain) {ドメイン：Order・DeliveryEstimate・業務規則};
    \node[accentbox,minimum width=92mm,minimum height=6mm,inner ysep=.5mm,below=of domain] (ports) {ポート：OrderRepository・DeliveryService・ResponseGenerator};
    \node[box,minimum width=92mm,minimum height=6mm,inner ysep=.5mm,below=of ports] (adapter2) {出力アダプタ：PostgreSQL／配送API／LLM API};
    \draw[flow] (adapter1)--(usecase);
    \draw[flow] (usecase)--(domain);
    \draw[flow] (domain)--(ports);
    \draw[flow] (adapter2)--(ports);
  \end{tikzpicture}
  \vspace{2mm}
\begingroup
\paragraph{内側の責務}
        注文と配送予定の事実、質問へ答える処理順序を表し、HTTP、SQL、特定ベンダーのSDKを直接知りません。

\par
\endgroup
\medskip
\begingroup
\paragraph{外側の責務}
        外部形式を内側の型へ変換し、注文の保存、配送照会、文章生成を具体技術で実装します。

\par
\endgroup
\par
\endgroup


\section{変更要求による評価}
\label{sec:04-software-engineering-21}

境界の良し悪しは、図の整い方ではなく、現実の変更を流したときに分かります。データベース、配送会社、入力経路、モデル提供者を一つずつ変え、業務上の事実まで不必要に揺れないかを確かめます。

\begingroup
  \noindent
  \begin{tabularx}{\textwidth}{@{}lYY@{}}
    \toprule
    変更要求 & 主に変える場所 & 変えずに保ちたい場所\\
    \midrule
    DBをPostgreSQLから別製品へ & Repositoryのアダプタ、設定、結合テスト & 注文規則、HTTP契約\\ \addlinespace[.7mm]
    配送会社を一社追加 & 配送アダプタ、配送先の選択規則 & 注文の保存、画面の契約\\ \addlinespace[.7mm]
    RESTへ音声入力を追加 & 入力アダプタ、入力の変換 & ユースケース、注文規則\\ \addlinespace[.7mm]
    LLM提供者を変更 & 生成アダプタ、出力検査、監視 & 配送予定の事実、注文規則\\
    \bottomrule
  \end{tabularx}
  \vspace{4mm}
\begingroup
\paragraph{波及の原因}
        内部型の漏出、共有状態、曖昧な失敗契約、循環依存がないかを確認します。

\par
\endgroup
\medskip
\begingroup
\paragraph{層の役割}
        小規模な処理へ同じ数の層を強制すると、間接参照だけが増えます。変更リスクに見合う境界を作ります。

\par
\endgroup
\par
\endgroup


\section{機械の外へ}
\label{sec:04-software-engineering-22}

一つのAIアプリは、注文の事実、配送サービス、HTTP、データベース、文章生成を抱えています。変化の理由ごとに責任を分け、インタフェースで約束し、モデルで協調を見渡し、テストで境界を確かめると、どこか一つの変更を局所へ閉じ込められます。オブジェクト指向や関数型は、その目的に応じて使う表現です。

設計図の中で、要求は入力アダプタからユースケースへ渡るところまで来ました。次に必要なのは、端末とサーバーの間を実際に越えることです。次章では箱と矢印で描いた「HTTP API」の内側へ入り、データが相手へ届くまでの物理的な旅を再開します。

\takeaway{良い設計は複雑さを消すのではなく、変わる理由ごとに置き場所を与え、変更の波が越える境界を少なくします。}


\section{旅をたどり直す}
\label{sec:04-software-engineering-23}

注文確認AIへ一つ変更が加わる場面を想像し、その波がどこまで届くかをたどります。次の問いでは手法名だけでなく、どの責任を守るために使うのかを説明してください。

\begingroup
  \begin{enumerate}
    \item ソフトウェアの規模が増えると、変更コストが増える理由を説明してください。
    \item 凝集度と結合度を、注文確認AIの例で説明してください。
    \item 情報隠蔽では、何をモジュール内へ隠しますか。
    \item 逐次・選択・反復が制御フローの理解を助ける理由を説明してください。
    \item クラス図とシーケンス図は、それぞれどの問いへ答えますか。
    \item 継承ではなく合成を選ぶべき関係を一つ示してください。
    \item 純粋関数と副作用を分けると、テストへどのような効果がありますか。
    \item ケースへLLM提供者の追加要求を流し、変更箇所を示してください。
  \end{enumerate}
\par
\endgroup
